\section*{Capítulo \ref{ch:b3:rings} --- Rings and Arithmetic}

\begin{solution}{exo:b3:rings:1}
(a) Evaluación $f \colon \Z[X] \to \Z[\iu]$, $P \mapsto P(\iu)$,
es un morfismo de anillo sobreyectivo ($a + bX \mapsto a + b\iu$). Núcleo:
dividir $P$ por \emph{mónico} $X^2 + 1$ en $\Z[X]$: $P = (X^2 +
1)Q + (bX + a)$con$a, b \in \Z$; entonces$P(\iu) = a + b\iu = 0$
sif $a = b = 0$. Entonces $\ker f = (X^2+1)$ y
\cref{thm:b3:rings:firstiso} concluye.

(b) El mismo cálculo con coeficientes $\R$: $\R[X]/(X^2+1)
\cong \C$--- este es el \emph{construcción} más limpio de$\C$.

(c) $X^2 + X + 1$ no tiene raíz en $\mathbb F_2$ ($0, 1 \mapsto 1$),
entonces, teniendo grado $2$, es \omterm{def:b3:rings:divisibility}{irreducible}: el cociente $\mathbb
F_4 = \mathbb F_2[X]/(X^2+X+1)$ es un campo
(\cref{prop:b3:rings:primemaximal}; $(P)$ \omterm{def:b3:rings:primemaximal}{máximo} en $K[X]$ cuando
$P$ es \omterm{def:b3:rings:divisibility}{irreducible}, ya que $K[X]$ es un \omterm{def:b3:rings:pidufd}{PID}: un \omterm{def:b3:rings:ideal}{ideal} ($(D) \supseteq
(P)$significa$D \mid P$). Sus cuatro elementos son$0, 1, \omega,
\omega + 1$donde$\omega = \bar X$, con$\omega^2 = \omega + 1$.
Tabla de multiplicar (elementos distintos de cero):
\[
\omega \cdot \omega = \omega + 1, \qquad
\omega(\omega + 1) = \omega^2 + \omega = 1, \qquad
(\omega+1)^2 = \omega^2 + 1 = \omega .
\]
Los elementos distintos de cero forman un grupo cíclico de orden $3$ generado por
$\omega$.
\end{solution}

\begin{solution}{exo:b3:rings:2}
(a) Sea $p$ \omterm{def:b3:rings:divisibility}{irreducible} en \omterm{def:b3:rings:pidufd}{UFD} y $p \mid ab$, digamos $ab =
pc$, con$a, b \neq 0$(de lo contrario, trivial). Si$a$o$b$ es una unidad,
$p$ divide al otro. De lo contrario, factorice $a$, $b$ y $c$ en
\omterm{def:b3:rings:divisibility}{irreductibles}: las dos factorizaciones de $ab$,
\[
(\text{factors of } a)(\text{factors of } b) = p \cdot
(\text{factors of } c),
\]
debe aceptar el pedido y \omterm{def:b3:rings:divisibility}{asociados}: $p$ es asociado de algunos
\omterm{def:b3:rings:divisibility}{irreducible} factor de $a$ o de $b$, por lo tanto lo divide.

(b) Sea $A$ un dominio finito y $x \neq 0$. La aplicación $y \mapsto
xy$es inyectivo ($xy = xy' \Rightarrow x(y - y') = 0 \Rightarrow
y = y'$), por lo tanto sobreyectivo ($A$finito):$1 = xy$para algunos$y$.

(c) Si $\mathfrak p$ es primo en un anillo finito $A$, entonces
$A/\mathfrak p$ es un dominio finito, por lo tanto, un campo por (b), por lo que
$\mathfrak p$ es \omterm{def:b3:rings:primemaximal}{máximo} (\cref{prop:b3:rings:primemaximal}).
\end{solution}

\begin{solution}{exo:b3:rings:3}
Normas: $N(3) = 9$, $N(1 \pm \iu\sqrt5) = 6$, $N(2 \pm \iu\sqrt5) =
9$. Las ecuaciones$x^2 + 5y^2 = 2$y$x^2 + 5y^2 = 3$ no tienen
soluciones enteras, por lo que ningún elemento tiene la norma $2$ o $3$. Un adecuado
La factorización de $3$ necesitaría dos factores de la norma $3$:
imposible --- $3$ es \omterm{def:b3:rings:divisibility}{irreducible}. Una factorización adecuada de $1
\pm \iu\sqrt5$(norma$6$) necesitaría factores de las normas$2, 3$:
imposible. Lo mismo para $2 \pm \iu\sqrt5$ (norma $9$: los factores
tiene norma $3$). ahora
\[
9 = 3 \cdot 3 = (2 + \iu\sqrt5)(2 - \iu\sqrt5),
\]
dos factorizaciones en \omterm{def:b3:rings:divisibility}{irreductibles}. son genuinamente
diferentes: las unidades son $\pm 1$ (norma $1$) y $2 \pm \iu\sqrt5
\neq \pm 3$. Entonces la unicidad falla, mientras que la existencia de
factorizaciones se mantienen en $\Z[\iu\sqrt5]$
(\cref{exo:b3:rings:10}(c)): la no factorialidad aquí es puramente una
falla de unicidad. (Consistentemente, \cref{prop:b3:rings:primeirred}:
estos \omterm{def:b3:rings:divisibility}{irreductibles} no son primos.)
\end{solution}

\begin{solution}{exo:b3:rings:4}
(a) Sean $a, b \in \Z[\iu]$, $b \neq 0$ y $a/b = x + \iu y \in
\C$. Elija números enteros$m, n$con$\abs{x - m} \leq \frac12$,
$\abs{y - n} \leq \frac12$ y establezca $q = m + \iu n$, $r = a - bq$.
entonces
\[
N(r) = N(b)\,\abs*{\tfrac ab - q}^2
\leq N(b)\Bigl(\tfrac14 + \tfrac14\Bigr) = \tfrac{N(b)}2 < N(b).
\]
Entonces $N$ es una función \omterm{def:b3:rings:pidufd}{euclidiano} ($N(r) < N(b)$ o $r = 0$).

(b) Si $uv = 1$ entonces $N(u)N(v) = 1$ con $N(u) \in \N$: $N(u) =
1$, es decir $x^2 + y^2 = 1$:$u \in \{\pm 1, \pm\iu\}$; por el contrario
estas son unidades.

(c) Para $\Z[\iu\sqrt2]$: el mismo redondeo da $\abs{a/b - q}^2
\leq \frac14 + \frac{2}4 = \frac34 < 1$: \omterm{def:b3:rings:pidufd}{euclidiano}; unidades:$x^2 +
2y^2 = 1$da$\pm 1$. Para$\Z[\iu\sqrt5]$ el límite se convierte en
$\frac14 + \frac54 = \frac32 > 1$: el argumento de redondeo falla ---
y debe fallar, ya que $\Z[\iu\sqrt5]$ ni siquiera es un \omterm{def:b3:rings:pidufd}{UFD}
(\cref{exo:b3:rings:3}), mientras que \omterm{def:b3:rings:pidufd}{euclidiano} implicaría \omterm{def:b3:rings:pidufd}{UFD}
(\cref{thm:b3:rings:euclideanpid,thm:b3:rings:pidufd}).
\end{solution}

\begin{solution}{exo:b3:rings:5}
(a) Si $x^n = 0$:
\[
(1 + x)\bigl(1 - x + x^2 - \dots + (-1)^{n-1}x^{n-1}\bigr)
= 1 + (-1)^{n-1}x^n = 1 .
\]
(b) En un dominio, $\deg(PQ) = \deg P + \deg Q$; $PQ = 1$ fuerzas
$\deg P = \deg Q = 0$ y $P, Q \in A^\times$: $A[X]^\times =
A^\times$. Sobre$\Z/4\Z$:$(1 + 2X)^2 = 1 + 4X + 4X^2 = 1$, por lo que$1
+ 2X$es una unidad de grado$1$(aquí$2$ es nilpotente; compare
(a)).

(c) $e^2 = e$ da $e(e - 1) = 0$, por lo que $e \in \{0, 1\}$ en un
dominio. Si $x^n = 0$ con $n \geq 1$ mínimo y $x \ne 0$, entonces
$n \geq 2$ y $x \cdot x^{n-1} = 0$ con ambos factores distintos de cero:
contradicción.
\end{solution}

\begin{solution}{exo:b3:rings:6}
(a) $K[X,Y]/(X,Y) \cong K$ (evaluar en $(0,0)$): un campo, por lo que
$(X,Y)$ es \omterm{def:b3:rings:primemaximal}{máximo}. Si $(X, Y) = (P)$: $P \mid X$ fuerzas (grados
en $Y$) $P \in K[X]$ y $P \mid Y$ luego fuerza a $P \in K$; $P =
0$es absurdo y$P \in K^\times$daría$(P) = K[X,Y]$,
propiedad contradictoria ($K[X,Y]/(X,Y) \cong K \neq 0$). entonces
$(X,Y)$ no es principal.

(b) La evaluación $P(X, Y) \mapsto P(X, X^2)$ asigna $K[X,Y]$ a
$K[X]$; su núcleo es $(Y - X^2)$: dividiendo por monic-in-$Y$
polinomio $Y - X^2$, $P = (Y - X^2)Q + R(X)$ y $P(X, X^2) =
R(X)$. Entonces$K[X,Y]/(Y - X^2) \cong K[X]$ --- el anillo de coordenadas de
una parábola, isomorfa a la de una recta.

La evaluación $P(X, Y) \mapsto P(X, X^{-1})$ asigna $K[X, Y]$ al
anillo $K[X, X^{-1}]$ de polinomios de Laurent. Su núcleo contiene
$(XY - 1)$; por el contrario, módulo $XY - 1$ cada clase tiene un
representante $R = \sum_{n \geq 0} a_nX^n + \sum_{m \geq 1} b_m
Y^m$(reemplace cada producto$XY$por$1$repetidamente) y$R(X,
X^{-1}) = \sum a_n X^n + \sum b_m X^{-m} = 0$fuerza todos los$a_n =
b_m = 0$. Por lo tanto$K[X, Y]/(XY - 1) \cong K[X, X^{-1}]$ --- el
Anillo de coordenadas de una hipérbola: la línea a la que se le quita un punto.

(c) $(Y - X^2)$ es primo (el cociente $K[X]$ es un dominio) pero no
\omterm{def:b3:rings:primemaximal}{máximo} ($K[X]$ no es un campo; concretamente $(Y - X^2) \subsetneq
(Y - X^2,\, X) \subsetneq K[X,Y]$).
\end{solution}

\begin{solution}{exo:b3:rings:7}
\emph{$X^5 - 12X^3 + 36X - 12$}: \omterm{thm:b3:rings:criteria}{Eisenstein} en $p = 3$ ($3 \mid
12, 36, 12$;$9 \nmid 12$;$3 \nmid 1$): \omterm{def:b3:rings:divisibility}{irreducible}. (En$p = 2$
\omterm{thm:b3:rings:criteria}{Eisenstein} falla: $4 \mid 12$.)

\emph{$X^4 + X + 1$}: reducir mod $2$. Sin raíz en $\mathbb F_2$;
la única cuadrática \omterm{def:b3:rings:divisibility}{irreducible} sobre $\mathbb F_2$ es $X^2 + X +
1$y$(X^2+X+1)^2 = X^4 + X^2 + 1 \neq X^4 + X + 1$. Entonces$X^4 +
X + 1$es \omterm{def:b3:rings:divisibility}{irreducible} sobre$\mathbb F_2$, por lo tanto, sobre$\Q$
(\cref{thm:b3:rings:criteria}(1); es mónico).

\emph{$X^4 + 4$}: reducible --- la identidad de Sophie Germain,
$X^4 + 4 = (X^2 - 2X + 2)(X^2 + 2X + 2)$.

\emph{$X^4 + 1$}: turno, $(X+1)^4 + 1 = X^4 + 4X^3 + 6X^2 + 4X +
2$: \omterm{thm:b3:rings:criteria}{Eisenstein} en$2$. Una factorización de$X^4+1$ cambiaría a
uno de $(X+1)^4 + 1$: \omterm{def:b3:rings:divisibility}{irreducible}.

\emph{$X^3 - X - 1$}: un cúbico es reducible sobre $\Q$ si tiene un
raíz racional; una raíz racional de un polinomio entero mónico es una
número entero que divide el término constante (teorema de la raíz racional: si
$(p/q)$ en términos más bajos es una raíz, $q \mid 1$, $p \mid -1$), y
$\pm 1$ no son raíces ($-1$ y $-1$): \omterm{def:b3:rings:divisibility}{irreducible}.
\end{solution}

\begin{solution}{exo:b3:rings:8}
(a) Para $\gcd(m, n) = 1$, \cref{thm:b3:rings:crt} suena
isomorfismo $\Z/mn\Z \cong \Z/m\Z \times \Z/n\Z$. Un elemento de un
el anillo producto es una unidad si ambas coordenadas lo son, entonces
$(\Z/mn\Z)^\times \cong (\Z/m\Z)^\times \times (\Z/n\Z)^\times$
y $\varphi(mn) = \varphi(m)\varphi(n)$. Para una potencia primaria, la
las no unidades de $\Z/p^k\Z$ son las clases de múltiplos de $p$:
$\varphi(p^k) = p^k - p^{k-1}$. Por lo tanto
\[
\varphi(n) = \prod_i \bigl(p_i^{\alpha_i} -
p_i^{\alpha_i-1}\bigr) = n \prod_{p \mid n}\Bigl(1 -
\frac1p\Bigr).
\]

(b) $M = 7 \cdot 11 \cdot 13 = 1001$. Idempotentes: $e_1 \equiv
(1, 0, 0)$:$143 = 11\cdot13 \equiv 3 \pmod 7$y$3 \cdot 5
\equiv 1$:$e_1 = 143 \cdot 5 = 715$.$e_2$:$91 \equiv 3
\pmod{11}$,$3 \cdot 4 \equiv 1$:$e_2 = 91\cdot4 = 364$.$e_3$:
$77 \equiv -1 \pmod{13}$: $e_3 = 77 \cdot 12 = 924$. entonces
\[
x \equiv 2\,e_1 + 5\,e_2 + 1\,e_3 = 1430 + 1820 + 924 = 4174
\equiv 170 \pmod{1001},
\]
y de hecho $170 = 24\cdot7 + 2 = 15\cdot11 + 5 = 13\cdot13 + 1$.
\end{solution}

\begin{solution}{exo:b3:rings:9}
(a) Si $x^n = 0$ y $y^m = 0$, la expansión binomial de
$(x+y)^{n+m}$ tiene todos los términos $x^iy^j$ con $i + j = n + m$, por lo que $i
\geq n$o$j \geq m$: cada término desaparece y$x + y$ es
nilpotente; $(ax)^n = a^nx^n = 0$: $\operatorname{Nil}(A)$ es un
\omterm{def:b3:rings:ideal}{ideal}. Si $\mathfrak p$ es primo y $x^n = 0 \in \mathfrak p$,
la inducción en $n$ da $x \in \mathfrak p$ ($x \cdot x^{n-1} \in
\mathfrak p$).

(b) Sean $a \notin \operatorname{Nil}(A)$ y $S = \{a^n : n \geq
1\}$, por lo que$0 \notin S$. El conjunto$\mathcal E$ de \omterm{def:b3:rings:ideal}{ideales} disjunto
de $S$ contiene $(0)$ y es inductivo (la unión de una cadena de
\omterm{def:b3:rings:ideal}{ideales} disjunto de $S$ es un \omterm{def:b3:rings:ideal}{ideal} disjunto de $S$): Zorn
proporciona $\mathfrak p \in \mathcal E$ \omterm{def:b3:rings:primemaximal}{máximo}. $\mathfrak p$ es
adecuado ($a \notin \mathfrak p$, desde $a \in S$). Primalidad: dejar
$x, y \notin \mathfrak p$. Por maximalidad, $\mathfrak p + (x)$ y
$\mathfrak p + (y)$ conoce a $S$: $a^m \in \mathfrak p + (x)$, $a^n
\in \mathfrak p + (y)$. Multiplicando,$a^{m+n} \in \mathfrak p +
(xy)$. Si$xy \in \mathfrak p$, entonces$a^{m+n} \in \mathfrak p \cap
S$: absurdo. Entonces$xy \notin \mathfrak p$ --- contrapositivo de
primalidad. Por tanto, todo elemento no nilpotente evita algún \omterm{def:b3:rings:primemaximal}{principal
ideal}; con (un),
$\operatorname{Nil}(A) = \bigcap_{\mathfrak p} \mathfrak p$.
\end{solution}

\begin{solution}{exo:b3:rings:10}
(a) La cadena $\ker f \subseteq \ker f^2 \subseteq \cdots$
estabiliza (\cref{prop:b3:rings:noethacc}): $\ker f^n = \ker
f^{n+1}$para algunos$n$. Deje$x \in \ker f$. Como$f$, por lo tanto$f^n$, es
sobreyectiva, $x = f^n(y)$ para algunos $y$; luego $f^{n+1}(y) = f(x) =
0$, entonces$y \in \ker f^{n+1} = \ker f^n$, es decir, $x = f^n(y) = 0$.

(b) $I_n = \{f \in \mathcal C(\intcc01, \R) : f\restriction_{
\intcc0{1/n}} = 0\}$es \omterm{def:b3:rings:ideal}{ideal} y$I_n \subseteq I_{n+1}$. el
la inclusión es estricta: $x \mapsto \max\bigl(0, x -
\frac1{n+1}\bigr)$desaparece en$\intcc0{\frac1{n+1}}$ pero no en
$\intcc0{\frac1n}$. Una cadena infinita estrictamente creciente
contradice \cref{prop:b3:rings:noethacc}.

(c) Suponga que el conjunto de no unidades distintas de cero que no admiten factorización
en \omterm{def:b3:rings:divisibility}{irreductibles} no está vacío. La familia correspondiente de \omterm{def:b3:rings:ideal}{ideales}
$\{(a)\}$ tiene un elemento \omterm{def:b3:rings:primemaximal}{máximo} $(a)$
(\cref{prop:b3:rings:noethacc}). El elemento $a$ no es
\omterm{def:b3:rings:divisibility}{irreducible} (un \omterm{def:b3:rings:divisibility}{irreducible} es su propia factorización), por lo que $a = bc$
con $b, c$ no unidades; $(a) \subseteq (b)$ es estricto (ya que $(a) =
(b)$daría$b = ad$,$a = adc$, entonces$dc = 1$:$c$ una unidad),
igualmente $(a) \subsetneq (c)$. Por maximidad, $b$ y $c$ ambos
factor en \omterm{def:b3:rings:divisibility}{irreductibles}; factores de concatenación $a$:
contradicción. Aplicado a $\Z[\iu\sqrt5]$ --- \omterm{def:b3:rings:noetherian}{noetheriano} como
cociente de $\Z[X]$ (\cref{thm:b3:rings:hilbert}, $\Z[\iu\sqrt5]
\cong \Z[X]/(X^2+5)$) --- esto muestra que existen factorizaciones allí;
\cref{exo:b3:rings:3} demostró que la unicidad es lo que falla.
\end{solution}

\begin{solution}{exo:b3:rings:11}
(a) \cref{exo:b3:rings:7}: turno y \omterm{thm:b3:rings:criteria}{Eisenstein} en $2$.

(b) Mod $2$: $X^4 + 1 = (X + 1)^4$. Ahora dejemos que $p$ sea impar. el
los cuadrados forman un subgrupo del índice $2$ en $(\Z/p\Z)^\times$: el
El morfismo $x \mapsto x^2$ tiene el núcleo $\{\pm 1\}$ (dos elementos: $X^2
- 1$tiene como máximo$2$raíces en un campo y$1 \neq -1$ para raíces impares).
$p$), por lo que su imagen tiene elementos $\frac{p-1}2$. En consecuencia, el
El producto de dos no cuadrados es un cuadrado (en el orden $2$ \omterm{thm:b3:groups:quotient}{cociente
grupo}, $\overline{xy} = \bar x\bar y$). Por lo tanto, \emph{al menos uno
de $-1$, $2$, $-2$ es un mod cuadrado $p$} (si $-1$ y $2$ no lo son,
$-2 = (-1)\cdot 2$ es). En cada caso $X^4 + 1$ factores mod $p$:
\begin{itemize}
  \item $-1 = c^2$: \quad $X^4 + 1 = X^4 - c^2 = (X^2 - c)(X^2 +
        c)$;
  \item $2 = c^2$: \quad $(X^2 + cX + 1)(X^2 - cX + 1) = X^4 + (2
        - c^2)X^2 + 1 = X^4 + 1$;
  \item $-2 = c^2$: \quad $(X^2 + cX - 1)(X^2 - cX - 1) = X^4 -
        (c^2 + 2)X^2 + 1 = X^4 + 1$.
\end{itemize}
Entonces $X^4+1$ es módulo reducible cada primo, pero \omterm{def:b3:rings:divisibility}{irreducible} supera
$\Q$: el criterio de reducción (\cref{thm:b3:rings:criteria}(1))
detecta irreductibilidad pero su fallo no prueba nada.

(Por razones estructurales: $p^2 - 1 = (p-1)(p+1)$ es un producto de
dos números pares consecutivos, entonces $8 \mid p^2 - 1$; el cíclico
grupo $\mathbb F_{p^2}^\times$ (ciclicidad demostrada en
\cref{ch:b3:galois}) contiene entonces un elemento $\zeta$ de orden
$8$, una raíz de $X^4 + 1$; su polinomio mínimo sobre $\mathbb
F_p$divide a$X^4+1$y tiene grado$\leq 2$---$X^4+1$ nunca puede
ser \omterm{def:b3:rings:divisibility}{irreducible} mod $p$.)
\end{solution}

\begin{solution}{exo:b3:rings:12}
(a) $(1-e)^2 = 1 - 2e + e^2 = 1 - e$. La aplicación $\varphi(x) =
(ex, (1-e)x)$ es aditivo y multiplicativo en el producto.
de los dos \omterm{def:b3:rings:ideal}{ideales}: $exey = e^2xy = e(xy)$, y $Ae$ es un
anillo conmutativo con unidad $e$ ($e\cdot ex = ex$). Inyectivo:
$ex = 0$ y $(1-e)x = 0$ suman $x = 0$. Sobreyectiva: $(ea,
(1-e)b)$es la imagen de$ea + (1-e)b$ (calcule ambas
componentes que utilizan $e(1-e) = 0$). Las unidades se asignan al estilo $(1, 0)$
pares correctamente: $\varphi(1) = (e, 1-e)$, la unidad del
producto.

(b) En un dominio, $e(e - 1) = 0$ fuerza a $e \in \{0, 1\}$: solo
Idempotentes triviales. En $\Z/12\Z$, resolviendo $e^2 \equiv e$: $e
\in \{0, 1, 4, 9\}$. El par no trivial$\{4, 9\}$:$4 + 9 =
13 \equiv 1$,$4\cdot9 = 36 \equiv 0$y$\Z/12\Z\cdot4 =
\{0, 4, 8\} \cong \Z/3\Z$(unidad$4$),$\Z/12\Z\cdot9 = \{0,
3, 6, 9\} \cong \Z/4\Z$(unidad$9$): la división CRT
$\Z/12\Z \cong \Z/4\Z\times\Z/3\Z$, con $9 \leftrightarrow
(1, 0)$y$4 \leftrightarrow (0, 1)$.

(c) Bajo el isomorfismo CRT $\Z/n\Z \cong \prod_i
\Z/p_i^{a_i}\Z$, el elemento$e_i$ con lo indicado
congruencias corresponde a la tupla con $1$ en el slot $i$ y
$0$ en otra parte: los idempotentes elementales del producto.
Por el contrario, una familia completa de idempotentes ortogonales.
($e_ie_j = 0$ para $i \neq j$, $\sum e_i = 1$) vuelve a montar el
descomposición del producto mediante (a), de forma inductiva. Los idempotentes son para
anillos qué son las proyecciones ortogonales a los espacios de Hilbert
(\cref{ch:b3:hilbert}): las coordenadas de un directo interno
descomposición.
\end{solution}

\begin{solution}{pb:b3:rings:1}
\textbf{1.} $N(z) = z\bar z$, entonces $N(zw) = zw\overline{zw} = z\bar
z\, w \bar w = N(z)N(w)$. Si$uv = 1$:$N(u)N(v) = 1$en$\N$, entonces
$N(u) = 1$, es decir\ $u \in \{\pm 1, \pm \iu\}$; los cuatro son unidades.
Brahmagupta: $(a^2+b^2)(c^2+d^2) = N\bigl((a + \iu b)(c + \iu
d)\bigr) = (ac - bd)^2 + (ad + bc)^2$.

\textbf{2.} Si $z = ab$, entonces $N(z) = N(a)N(b)$ es primo, entonces
$N(a) = 1$ o $N(b) = 1$: un factor es una unidad. Como $N(z) > 1$,
$z$ no es cero ni una unidad: \omterm{def:b3:rings:divisibility}{irreducible} --- y primo, ya que
$\Z[\iu]$ es un \omterm{def:b3:rings:pidufd}{UFD}
(\cref{thm:b3:rings:euclideanpid,thm:b3:rings:pidufd,%
lem:b3:rings:bezout}).

\textbf{3.} $\pi$ divide $N(\pi) = \pi\bar\pi \geq 2$, un
entero; factorizar $N(\pi)$ en números primos y usarlo
$\pi$ es primo, $\pi \mid p$ para algún número primo $p$. si también
$\pi \mid q \neq p$: Bézout en $\Z$ da $1 = up + vq$, entonces
$\pi \mid 1$ --- absurdo: $p$ es único. De $p = \pi\gamma$:
$p^2 = N(p) = N(\pi)N(\gamma)$ con $N(\pi) \neq 1$, entonces $N(\pi)
\in \{p, p^2\}$.

\textbf{4.} Sea $\pi$ un primo gaussiano que divide $p$, $p =
\pi\gamma$. Si$N(\pi) = p^2$:$N(\gamma) = 1$, entonces$p$ es un
asociado de $\pi$, en sí mismo un primo gaussiano; y no gaussiano
Prime tiene la norma $p$ (si $N(\rho) = p$ entonces $\rho \mid
\rho\bar\rho = p$, y$p$Prime en$\Z[\iu]$ forzaría
$\rho$ asociado a $p$, dando $N(\rho) = N(p) = p^2 \neq
p$).
Si $N(\pi) = p$: escribir $\pi = a + \iu b$, $p = \pi\bar\pi = a^2
+ b^2$.

\textbf{5.} En el grupo abeliano $(\Z/p\Z)^\times$, empareje cada uno
elemento con su inversa. Los elementos autoinversos son las raíces.
de $X^2 - 1$: exactamente $\pm 1$ (como máximo dos raíces en un campo).
El producto de todos los elementos es entonces $1 \cdot (-1) \cdot \prod
(\text{pairs } k k^{-1}) = -1$:$(p-1)! \equiv -1 \pmod p$. (Para
$p = 2$: $1! \equiv -1$.)

\textbf{6.} Escriba $(p-1)! = \prod_{k=1}^m k \cdot
\prod_{k=m+1}^{p-1}k$con$m = \frac{p-1}2$. en el segundo
sustituto del producto $k = p - j$, $j = 1, \dots, m$: módulo $p$,
$\prod_{j=1}^m (p - j) \equiv (-1)^m m!$. Por lo tanto, $-1 \equiv
(-1)^m (m!)^2$, es decir, $(m!)^2 \equiv (-1)^{m+1} \pmod p$.

\textbf{7.} Si $p \equiv 1 \pmod 4$, $m$ es par y pregunta 6
da $(m!)^2 \equiv -1$: una raíz cuadrada de $-1$. Por el contrario, si
$x^2 \equiv -1 \pmod p$ ($p$ impar), luego $x^4 = 1 \neq x^2$: $x$
tiene orden $4$ en $(\Z/p\Z)^\times$, por lo que $4 \mid p - 1$ (Lagrange).
Y $p = 2$: $1^2 = 1 \equiv -1$. Conclusión: $-1$ es un cuadrado
mod $p$ si $p = 2$ o $p \equiv 1 \pmod 4$.

\textbf{8.} Con $x^2 \equiv -1$: $p \mid x^2 + 1 = (x + \iu)(x -
\iu)$. Si$p$ fuera un primo gaussiano, dividiría un factor; pero
$\frac xp \pm \frac \iu p \notin \Z[\iu]$. Entonces $p$ no es un
prima gaussiana; por la dicotomía (pregunta 4) --- $p$ no es primo
significa la segunda rama --- $p = a^2 + b^2$.

\textbf{9.} Los cuadrados son $\equiv 0$ o $1 \pmod 4$, por lo que $a^2 + b^2
\in \{0, 1, 2\} \pmod 4$: un primo$p \equiv 3 \pmod 4$ no es un
suma de dos cuadrados. Según la pregunta 4, la rama $N(\pi) = p$ ($p =
a^2+b^2$) es imposible:$p$ sigue siendo un primo gaussiano.

\textbf{10.} $(1 + \iu)^2 = 2\iu$, entonces $2 = -\iu(1 + \iu)^2$; y
$N(1 + \iu) = 2$ es primo, por lo que $1 + \iu$ es un primo gaussiano
(pregunta 2).

\textbf{11.} Cada primo gaussiano divide exactamente a un primo
número $p$ (pregunta 3); listado por casos: $p = 2$ da el
\omterm{def:b3:rings:divisibility}{asociados} de $1 + \iu$; $p \equiv 3 \pmod 4$ se entrega a $p$
(pregunta 9); $p \equiv 1 \pmod 4$ entrega la pareja $\pi, \bar\pi$
de la norma $p$ (preguntas 4 y 8). El par es genuino: $\bar\pi \in
\{\pm\pi, \pm\iu\pi\}$forzaría, escribiendo$\pi = a + \iu b$,
ya sea $b = 0$, $a = 0$ o $a = \pm b$, lo que da $p = a^2 + b^2 \in
\{a^2, 2a^2\}$ --- imposible para un número primo impar. Comprobar:
$5 = (2 + \iu)(2 - \iu)$, $N(2\pm\iu) = 5$; $3$: prima de la norma
$9$.

\textbf{12.} Escriba $n = 2^{\alpha}\prod_i p_i^{\beta_i} \prod_j
q_j^{2\gamma_j}$con$p_i \equiv 1$,$q_j \equiv 3 \pmod 4$.
Cada factor es una suma de dos cuadrados: $2 = 1^2 + 1^2$; $p_i = a^2
+ b^2$(pregunta 8);$q_j^{2\gamma_j} = (q_j^{\gamma_j})^2 +
0^2$. La identidad Brahmagupta (pregunta 1) propaga la
propiedad al producto $n$.

\textbf{13.} Dejemos $n = a^2 + b^2 = N(a + \iu b)$ y $q \equiv 3
\pmod 4$,$q \mid n$. El primo gaussiano$q$ (pregunta 9) divide
$(a + \iu b)(a - \iu b)$, de ahí uno de los dos factores --- digamos $q
\mid a + \iu b$(el otro caso es idéntico). Pero luego$\frac{a +
\iu b}{q} = \frac aq + \iu \frac bq \in \Z[\iu]$se lee como$q
\mid a$\emph{y}$q \mid b$en$\Z$. Por lo tanto$q^2 \mid n$ y
$\frac n{q^2} =
\bigl(\frac aq\bigr)^2 + \bigl(\frac bq\bigr)^2$. por fuerte
inducción en $n$, el exponente de $q$ en $n/q^2$ es par; el de
$n$ también lo es.

\textbf{14.} \emph{Teorema (Fermat).} Un número entero positivo es una suma
de dos cuadrados si y sólo si todo primo $\equiv 3 \pmod 4$
ocurre en él con un exponente par. ---$2025 = 3^4 \cdot 5^2$:
exponente de $3$ par, sí ($2025 = 45^2 + 0^2 = 27^2 + 36^2$).
$2026 = 2 \cdot 1013$ con $1013 \equiv 1 \pmod 4$ principal: sí
($1013 = 22^2 + 23^2$ y Brahmagupta con $2 = 1^2+1^2$:
$2026 = (22 - 23)^2 + (22 + 23)^2 = 1^2 + 45^2$). $2027$ es un
principal $\equiv 3 \pmod 4$: no.

\textbf{15.} Deje $p = a^2 + b^2 = c^2 + d^2$ con positivo
enteros, $p \equiv 1 \pmod 4$ y $\pi$ un primo gaussiano con
$p = \pi\bar\pi$ (pregunta 4). Tanto $a + \iu b$ como $c + \iu d$
tienen norma $p$, por lo tanto, los primos gaussianos (pregunta 2) se dividen
$p = (a+\iu b)(a - \iu b)$; por unicidad de factorización, $c +
\iu d$es asociado de$a + \iu b$o de$a - \iu b$:
\[
c + \iu d \in \{\pm(a \pm \iu b),\ \pm\iu(a \pm \iu b)\}
= \{\pm a \pm \iu b,\ \pm b \pm \iu a\}.
\]
La positividad de $c, d$ deja a $c + \iu d \in \{a + \iu b, b + \iu
a\}$:$\{c, d\} = \{a, b\}$.

\textbf{16.} Si $n = a^2 - b^2 = (a-b)(a+b)$: los dos factores
tienen la misma paridad, por lo que $n$ es impar (ambos impares) o divisible por
$4$ (ambos pares) --- nunca $\equiv 2 \pmod 4$. Por el contrario, $n$
impar: $n = \bigl(\frac{n+1}2\bigr)^2 -
\bigl(\frac{n-1}2\bigr)^2$;$n = 4m$:$n = (m+1)^2 - (m-1)^2$.
La respuesta es una simple condición de congruencia, con una línea
identidad detrás de esto: las diferencias de cuadrados no conllevan aritmética
profundidad, y el contraste con las sumas es el objetivo de este
problema.

\textbf{17.} $(a, b) \mapsto z = a + \iu b$ es una biyección
entre representaciones y $\{z : N(z) = n\}$. Factorizar $z$ en
el \omterm{def:b3:rings:pidufd}{UFD} $\Z[\iu]$ usando la clasificación (pregunta 11): hasta
una unidad, $z = (1+\iu)^{a'}\prod_j\pi_j^{s_j}\bar\pi_j^{t_j}
\prod_kq_k^{u_k}$, y tomando normas ($N(1+\iu) = 2$,
$N(\pi_j) = N(\bar\pi_j) = p_j$, $N(q_k) = q_k^2$):
\[
n = 2^{a'}\prod_jp_j^{s_j + t_j}\prod_kq_k^{2u_k} .
\]
Exponentes coincidentes: $a' = a$, $s_j + t_j = b_j$, $2u_k = c_k$
--- \omterm{def:b3:groups:derived}{soluble} si cada $c_k$ es par, y entonces $u_k = c_k/2$ es
forzado mientras $s_j \in \intint0{b_j}$ es libre. Datos distintos
$(u, (s_j))$ proporciona $z$ no asociados con \emph{mismo}
norma; la unidad $u \in \{\pm1, \pm\iu\}$ (4 opciones) y luego
enumera cada clase asociada sin repetición (dos iguales
los productos violarían la unicidad de la factorización ---
$\pi_j$ y $\bar\pi_j$ no están asociados ya que $p_j =
\pi_j\bar\pi_j$no está ramificado). Total:$r_2(n) =
4\prod_j(b_j + 1)$y$0$si algún$c_k$ es impar.

\textbf{18.} $\chi(dd') = \chi(d)\chi(d')$ está marcado mod $4$
(impar $\times$ impar cubre los cuatro casos de signos; cualquier cosa par
da $0 = 0$). Para coprime $m, n$, los divisores de $mn$ son
únicamente $d = d_1d_2$ con $d_1 \mid m$, $d_2 \mid n$:
$\sum_{d \mid mn}\chi(d) = \bigl(\sum_{d_1\mid
m}\chi(d_1)\bigr)\bigl(\sum_{d_2\mid n}\chi(d_2)\bigr)$:
multiplicativo. Poderes primos: en $2^a$, solo $d = 1$ es impar:
suma $= 1$. En $p^b$ con $p \equiv 1$: todos $\chi(p^i) = 1$,
suma $= b + 1$. En $q^c$ con $q \equiv 3$: $\chi(q^i) =
(-1)^i$, suma alterna$= 1$($c$par) o$0$($c$ impar).

\textbf{19.} Las dos funciones multiplicativas
$\frac14r_2$ (pregunta 17) y $\sum_{d\mid n}\chi(d)$
(pregunta 18) está de acuerdo con todos los poderes primarios --- $1$ en $2^a$;
$b + 1$ en $p^b$; $\mathbf 1_{c\ \mathrm{even}}$ en $q^c$ ---
de ahí que coincida en todas partes: la fórmula de Jacobi, con $\sum_{d\mid
n}\chi(d) = d_1(n) - d_3(n)$ ordenando divisores. Cheques:
$r_2(3) = 0 = 4(1 - 1)$; $r_2(5) = 8 = 4(2 - 0)$
($(\pm1,\pm2), (\pm2,\pm1)$); $r_2(9) = 4 = 4(2 - 1)$
(divisores $1, 9 \equiv 1$; $3 \equiv 3$; representaciones
$(\pm3, 0), (0, \pm3)$); $r_2(25) = 12 = 4(3 - 0)$. Para $65 =
5\cdot13$:$r_2 = 4\cdot2\cdot2 = 16$, de$65 = 1 + 64 = 16
+ 49$: los dieciséis pares$(\pm1, \pm8), (\pm8, \pm1), (\pm4,
\pm7), (\pm7, \pm4)$.

\textbf{20.} $\sum_{n \leq x}r_2(n)$ cuenta los pares $(a, b)$
con $0 < a^2 + b^2 \leq x$, es decir,\ los puntos de la red del
disco cerrado $D_{\sqrt x}$ menos el origen. asignar a cada uno
punto de red $P$ el cuadrado unitario $P + \intco01^2$: estos
los cuadrados recubren el avión. Todo cuadrado unido a un punto de
$D_{\sqrt x}$ se encuentra en $D_{\sqrt x + \sqrt2}$, y cada cuadrado
reunión $D_{\sqrt x - \sqrt 2}$ está adjunta a un punto de
$D_{\sqrt x}$ (el cuadrado tiene diámetro $\sqrt 2$): comparando
áreas,
\[
\pi(\sqrt x - \sqrt2)^2 \leq \#\{\text{lattice points in }
D_{\sqrt x}\} \leq \pi(\sqrt x + \sqrt 2)^2,
\]
y ambos límites son $\pi x + O(\sqrt x)$. Restando el
El origen no cambia nada con esta precisión.

\textbf{21.} Por Jacobi (pregunta 19) e intercambiando el pedido
de sumatoria ($n = dm$):
\[
\frac14\sum_{n\leq x}r_2(n) = \sum_{n \leq x}\sum_{d \mid
n}\chi(d) = \sum_{d \leq x}\chi(d)\,\#\{m : dm \leq x\}
= \sum_{d\leq x}\chi(d)\Bigl\lfloor\frac xd\Bigr\rfloor,
\]
que es $\frac{\pi x}4 + O(\sqrt x)$ según la pregunta 20. Eliminar
los pisos: $\lfloor x/d\rfloor = x/d + O(1)$, pero sumando
$O(1)$ sobre $d \leq x$ es demasiado burdo; en su lugar use que el
las sumas parciales de $\chi$ están acotadas ($0, 1, 1, 0$ cíclicamente),
así por Abel resumen $\sum_{d\leq x}\chi(d)\{x/d\}$, cuyo
términos que agrupamos en pares $d \equiv 1, 3$, es $O(\sqrt x)$ ---
alternativamente y más simplemente: dividir en $\sqrt x$. Para $d \leq
\sqrt x$, reemplace$\lfloor x/d\rfloor$por$x/d + O(1)$: error
$O(\sqrt x)$. Para $d > \sqrt x$, $\lfloor x/d\rfloor$ toma
cada valor $v < \sqrt x$ en un intervalo de $d$ consecutivos,
en el que la suma $\chi$ es $O(1)$: error total $O(\sqrt x)$ por
sumando los valores $\leq \sqrt x$ de $v$, mientras que
$\sum_{d > \sqrt x}\chi(d)\frac xd = O(\sqrt x)$ por
colas de series alternas ($x\sum_{d>\sqrt x}\chi(d)/d =
x\,O(1/\sqrt x)$). Por lo tanto
\[
x\sum_{d \leq x}\frac{\chi(d)}d = \frac{\pi x}4 + O(\sqrt x),
\qquad\text{i.e.}\qquad
\sum_{d\leq x}\frac{\chi(d)}d = \frac\pi4 +
O\Bigl(\frac1{\sqrt x}\Bigr),
\]
y dejando $x \to \infty$: $1 - \frac13 + \frac15 - \dots =
\frac\pi4$.

\textbf{22.} Si $n \equiv 3 \pmod4$ fuera $a^2 + b^2$: cuadrados
son $\equiv 0, 1 \pmod 4$ y $a^2 + b^2 \in \{0, 1, 2\}$
mod $4$ --- imposible. (El criterio de la pregunta 17 dice que
igual: $n \equiv 3 \pmod 4$ obliga a algunos números primos $\equiv 3$ a ser impares
exponente.) Entonces, los números enteros representables evitan un residuo completo
clase: densidad $\leq \frac34$. El promedio $\pi$ de $r_2$
se concentra en unos pocos números enteros: $n = \prod_{j\leq k}p_j$
(primos distintos $\equiv 1 \bmod 4$) tiene $r_2(n) = 4\cdot2^k$
representaciones --- ilimitadamente muchas --- por lo que un conjunto disperso de
$n$ puede llevar el promedio completo, exactamente como el promedio de una lotería
La recompensa coexiste con una pérdida casi segura. Landau's
$\#\{n \leq x \text{ representable}\} \sim
Cx/\sqrt{\log x}$lo confirma: densidad$0$, media$\pi$.

\textbf{23.} \emph{Necesidad.} Deje $n = a^2 + b^2$ con
$\gcd(a, b) = 1$. Si un primo $q \equiv 3 \pmod 4$ divide $n$,
La pregunta 13 muestra $q \mid a$ y $q \mid b$: contradicción. si
$4 \mid n$: los cuadrados son $\equiv 0, 1 \pmod 4$, por lo que $a^2 + b^2
\equiv 0 \pmod 4$fuerza a$a^2 \equiv b^2 \equiv 0$, es decir $a,
b$ambos pares: contradicción. \emph{Suficiencia.} Escribir$n =
2^{\alpha}\prod_jp_j^{b_j}$con$\alpha \leq 1$y$p_j
\equiv 1 \pmod 4$, y configurar$z = (1+\iu)^{\alpha}\prod_j
\pi_j^{b_j} = a + \iu b$, de norma$n$. Supongamos un primo$t$
divide $\gcd(a, b)$; luego $t \mid z$ en $\Z[\iu]$. Si $t
\equiv 3 \pmod 4$:$t \mid N(z) = n$, excluido. Si$t \equiv 1
\pmod 4$:$t = \pi_t\bar\pi_t$, entonces$\bar\pi_t \mid z$; pero el
La factorización de $z$ no contiene primos conjugados ($\pi_j$ y
$\bar\pi_j$ son no asociados, pregunta 17), contradiciéndose
factorización única. Si $t = 2 = -\iu(1+\iu)^2$: entonces
$(1+\iu)^2 \mid z$, forzando $\alpha \geq 2$, excluido. Por lo tanto
$\gcd(a, b) = 1$: la representación es primitiva.

\textbf{24.} $a$ es impar ($\gcd(a, b) = 1$, $b$ par), por lo que $c^2
= a^2 + b^2$es impar y$c$es impar. Sea$\delta$ un común
Divisor primo gaussiano de $a + \iu b$ y $a - \iu b$:
divide su suma $2a$ y su diferencia $2\iu b$, por lo tanto
$2a$ y $2b$; una relación Bézout $ua + vb = 1$ luego da
$\delta \mid 2$, por lo que $\delta$ está asociado a $1 + \iu$ y
$N(\delta) = 2$ divide $N(a + \iu b) = c^2$, lo cual es impar:
contradicción. Entonces $a + \iu b$ y $a - \iu b$ son coprimos con
producto $c^2$; en el \omterm{def:b3:rings:pidufd}{UFD} $\Z[\iu]$, cada primo gaussiano de
$c^2$ ocurre con un exponente par y se divide completamente en uno
de los dos factores coprimos, de donde $a + \iu b = u(m + \iu
n)^2 = u\bigl(m^2 - n^2 + 2\iu mn\bigr)$con$u$ una unidad. el
Las opciones $u = \pm\iu$ hacen que la parte real $\mp 2mn$ sea pareja ---
imposible, $a$ es impar. Las opciones $u = \pm1$ dan, después
ajustando los signos de $m, n$ e intercambiando sus nombres para hacer
todo positivo, $a = m^2 - n^2$, $b = 2mn$ con $m > n
\geq 1$; y$c^2 = N(m + \iu n)^2$da$c = m^2 + n^2$. un
divisor común de $m$ y $n$ dividiría $a$ y $b$: $\gcd
(m, n) = 1$; y$m \equiv n \pmod 2$haría que$a$ fuera par:
paridades opuestas. Comprobaciones: $(m, n) = (2, 1)$ da $(3, 4,
5)$;$(m, n) = (5, 2)$proporciona$(25 - 4, 20, 25 + 4) = (21, 20,
29)$y$441 + 400 = 841 = 29^2$.

\textbf{25.} La fórmula de Jacobi $r_2(n) = 4(d_1(n) - d_3(n))$
da, para $n = 1, \dots, 25$:
\[
4,\ 4,\ 0,\ 4,\ 8,\ 0,\ 0,\ 4,\ 4,\ 8,\ 0,\ 0,\ 8,\ 0,\ 0,\
4,\ 8,\ 4,\ 0,\ 8,\ 0,\ 0,\ 0,\ 0,\ 12,
\]
distinto de cero exactamente en $n = 1, 2, 4, 5, 8, 9, 10, 13, 16, 17, 18,
20, 25$(por ejemplo$r_2(15) = 0$: divisores$1, 5 \equiv 1$
y saldo $3, 15 \equiv 3$; $r_2(20) = 8$: divisores $1, 5
\equiv 1$, ninguno$\equiv 3$). El total es$4 + 4 + 4 + 8 + 4 +
4 + 8 + 8 + 4 + 8 + 4 + 8 + 12 = 80$. El lado divisor: el
impar $d \leq 25$ contribuye
\[
25 - 8 + 5 - 3 + 2 - 2 + 1 - 1 + 1 - 1 + 1 - 1 + 1 = 20,
\]
leyendo $\chi(d)\lfloor 25/d\rfloor$ para $d = 1, 3, 5, \dots,
25$; y$4 \cdot 20 = 80$, como lo predicen las preguntas 21
identidad. Puntos de red del disco cerrado de radio $5$:
los puntos $80$ con $1 \leq a^2 + b^2 \leq 25$ más el
origen, es decir\ $81$; y $\pi x = 25\pi \approx 78.54$, un
error de aproximadamente $2.46$, cómodamente dentro del $O(\sqrt x)$
banda de la pregunta 20 ($\sqrt x = 5$).
\end{solution}